Quantum circuits require validation and testing to ensure that they perform as expected. 
Classical simulation is a powerful tool for this purpose, but it is limited by the exponential growth of complexity with the number of qubits.
Nevertheless, using the ZX-calculus, a graphical language for quantum mechanics,
we can attempt to reduce how exponential this growth is.
Recent developments have allowed focus to be placed on complexity being dependent on the number of $T$-gates that are in the circuit.
Given $t$ $T$-gates, if the complexity of classically simulating the circuit is $O(2^{\alpha t})$,
the aim is to reduce the value of $\alpha$ as much as possible, where $0<\alpha<1$.
This thesis presents a heuristic approach toward exploiting the structure of the
ZX-diagrams used to represent the quantum circuits, in order to reduce $\alpha$.
We conceptualise a formal definition of the heuristic method, applying it to existing approaches,
and also develop new heuristics that take advantage of the structure made available in the ZX-diagrams.
These heuristics perform well in practice, 
boosting the number of $T$-gates that can be simulated classically by up to three times as many
on certain classes of circuits.
